\chapter{Polar Coordinates}
\label{ch:ch14-polar-coordinates}

\section{A Different Way to Locate a Point}

Throughout this book, points in the plane have been described using
Cartesian coordinates. A point is located by measuring a horizontal
displacement and a vertical displacement from the origin, producing an
ordered pair

\[
(x,y).
\]

This system is natural whenever horizontal and vertical directions play a
special role, and it has been used repeatedly since Chapter~2. Yet many
geometric situations are concerned not with horizontal and vertical
motion, but with distance and direction. A navigator determining the
location of a ship may record a bearing and a range; a radar station may
measure an angle and a distance; an astronomer may specify the direction
and apparent separation of an object from a fixed observation point.

Trigonometry has spent much of this book developing exactly these two
quantities. Distances were introduced through the geometry of triangles
and vectors, while angles became increasingly important through rotation,
circular motion, and periodic functions. It is therefore natural to ask
what happens if position is described directly in terms of distance and
angle rather than horizontal and vertical displacement.

The resulting coordinate system is called the \emph{polar coordinate
system}. Instead of measuring how far a point lies along two perpendicular
axes, it measures how far the point lies from the origin and in what
direction it is found.

\section{Polar Coordinates}

In the polar coordinate system, the origin is called the \emph{pole}, and
the positive $x$-axis serves as the \emph{polar axis}. A point is located
by two numbers:

\[
(r,\theta),
\]

where $r$ is the distance from the pole and $\theta$ is the angle measured
from the positive $x$-axis.

\begin{definition}
A point in \emph{polar coordinates} is represented by an ordered pair
$(r,\theta)$, where $r$ is the distance from the origin and $\theta$ is
the angle measured from the positive $x$-axis.
\end{definition}

The angle convention is exactly the one established earlier for the unit
circle. Positive angles are measured counterclockwise and negative angles
clockwise. No new notion of angle is being introduced; the polar
coordinate system simply uses the rotational framework already developed
throughout the book.

For example, $(4,\pi/6)$ represents the point lying four units from the
origin in the direction determined by an angle of $\pi/6$ radians.
Likewise, $(2,3\pi/4)$ represents the point two units from the origin at
an angle of $135^\circ$. The distance and direction are given directly
rather than being inferred from rectangular coordinates.

\section{Converting Polar Coordinates to Rectangular Coordinates}

The relationship between polar and Cartesian coordinates follows directly
from the trigonometry of the unit circle. Consider a point with polar
coordinates $(r,\theta)$. Dropping a perpendicular from the point to the
$x$-axis forms a right triangle whose hypotenuse has length $r$ and whose
acute angle is $\theta$. From the definitions of sine and cosine,
$\cos\theta = x/r$ and $\sin\theta = y/r$, and solving for $x$ and $y$
gives the conversion formulas.

\begin{theorem}
If a point has polar coordinates $(r,\theta)$, then its rectangular
coordinates are $x = r\cos\theta$ and $y = r\sin\theta$.
\end{theorem}

The formulas should appear familiar. They are exactly the coordinate
equations of the unit circle, except that the radius is no longer fixed at
one. Every point is obtained by taking the unit-circle coordinates
associated with $\theta$ and scaling them by the distance $r$.

\begin{example}
Convert $\left(6,\dfrac{\pi}{3}\right)$ to rectangular coordinates. Using
the conversion formulas,
\[
x = 6\cos\frac{\pi}{3} = 6\left(\frac12\right) = 3, \qquad
y = 6\sin\frac{\pi}{3} = 6\left(\frac{\sqrt3}{2}\right) = 3\sqrt3.
\]
Therefore the rectangular coordinates are $(3, 3\sqrt3)$.
\end{example}

\section{Converting Rectangular Coordinates to Polar Coordinates}

The conversion from polar coordinates to rectangular coordinates follows
directly from the definitions of sine and cosine. The reverse conversion
requires a little more care. Suppose a point is given in rectangular
coordinates as $(x,y)$. The distance from the origin is immediately
determined by the distance formula, since the point, the origin, and the
coordinate axes form a right triangle: $r^2 = x^2+y^2$, so
\[
r = \sqrt{x^2+y^2}.
\]
This should look familiar. It is exactly the magnitude formula for
vectors introduced in Chapter~2; the radial coordinate $r$ is simply the
length of the position vector from the origin to the point.

To determine the angle, observe that $\tan\theta = y/x$, provided
$x\neq0$, which suggests $\theta = \arctan(y/x)$. This expression is not
always correct. The difficulty arises because the tangent function does
not uniquely determine a quadrant: the points $(1,1)$ and $(-1,-1)$ both
satisfy $y/x = 1$, yet they lie in entirely different directions from the
origin. The inverse tangent function introduced in Chapter~7 returns
values only in its principal range, $-\pi/2 < \theta < \pi/2$, which
covers Quadrants~I and IV but excludes Quadrants~II and III. The value
returned by $\arctan(y/x)$ must therefore be interpreted in light of the
point's actual location.

\begin{theorem}
If a point has rectangular coordinates $(x,y)$, then $r = \sqrt{x^2+y^2}$,
and the angle $\theta$ is determined by $\tan\theta = y/x$ together with
the requirement that $\theta$ lie in the correct quadrant.
\end{theorem}

The quadrant check is not a technical detail but an essential part of the
conversion process; forgetting it often produces a direction that differs
from the correct answer by $\pi$ radians.

\begin{example}
Convert the point $(-3,3)$ to polar coordinates. First, $r =
\sqrt{(-3)^2+3^2} = \sqrt{18} = 3\sqrt2$. Next, $\tan\theta = 3/(-3) =
-1$, and $\arctan(-1) = -\pi/4$. The point, however, lies in
Quadrant~II, not Quadrant~IV, so the correct angle is $\theta = 3\pi/4$.
Hence the polar coordinates are $\left(3\sqrt2, \dfrac{3\pi}{4}\right)$.
\end{example}

\section{Non-Uniqueness of Polar Coordinates}

One of the most surprising features of polar coordinates is that they do
not provide a unique representation of a point. In the rectangular
system, a point corresponds to exactly one ordered pair: $(3,4)$ cannot
also be represented by some different rectangular coordinate pair, and
the relationship between points and coordinates is one-to-one. Polar
coordinates behave differently. Consider the point $\left(2,
\dfrac{\pi}{3}\right)$; if one full revolution is added to the angle, the
direction does not change, so $\left(2,\dfrac{\pi}{3}\right) =
\left(2,\dfrac{\pi}{3}+2\pi\right)$, and adding further revolutions
produces infinitely many coordinate pairs representing the same point.

\begin{theorem}
For any integer $k$, $(r,\theta) = (r,\theta+2\pi k)$.
\end{theorem}

This result is a direct consequence of the periodicity of sine and
cosine; the coordinate system inherits the same rotational symmetry
already encountered in the study of coterminal angles in Chapter~4. An
even more surprising phenomenon occurs when negative values of $r$ are
allowed. The point $(r,\theta)$ lies exactly opposite the point
$(r,\theta+\pi)$ on the same line through the origin, so changing the
sign of the radius reverses that direction, bringing the point back to
its original location, and consequently $(r,\theta) = (-r,\theta+\pi)$.

\begin{example}
The point $\left(3,\dfrac{\pi}{6}\right)$ can also be written as
$\left(-3,\dfrac{7\pi}{6}\right)$; both descriptions represent the same
point in the plane.
\end{example}

The non-uniqueness of polar coordinates is not a defect of the system; it
is a consequence of describing position through distance and rotation.
Rotations naturally repeat every full revolution, and the coordinate
system reflects that periodicity directly.

\begin{algorithmproc}{How to Convert Between Rectangular and Polar Coordinates}
To convert from polar coordinates $(r,\theta)$ to rectangular
coordinates, use $x = r\cos\theta$ and $y = r\sin\theta$. To convert from
rectangular coordinates $(x,y)$ to polar coordinates, first compute $r =
\sqrt{x^2+y^2}$, then compute $\tan\theta = y/x$ and determine the
appropriate quadrant from the signs of $x$ and $y$, and finally choose an
angle $\theta$ in that quadrant satisfying the tangent equation.
\end{algorithmproc}

\begin{example}
Convert the point $(-4,-4)$ to polar coordinates. First, $r =
\sqrt{(-4)^2+(-4)^2} = \sqrt{32} = 4\sqrt2$. Next, $\tan\theta =
(-4)/(-4) = 1$, and $\arctan(1) = \pi/4$; the point, however, lies in
Quadrant~III, whose corresponding angle is $\theta = 5\pi/4$. Therefore
$(-4,-4) = \left(4\sqrt2, \dfrac{5\pi}{4}\right)$.
\end{example}

\begin{practicenow}
Convert the point $(2,-2\sqrt3)$ to polar coordinates, choosing an angle
in the correct quadrant.
\end{practicenow}

\section{Graphing in Polar Coordinates}

Changing coordinates does more than alter notation; it changes which
geometric objects appear simple. In Cartesian coordinates, a horizontal
line is described by an equation as simple as $y=3$, precisely because
the coordinate system itself is built from horizontal and vertical
directions. Polar coordinates instead privilege distance and angle, so
objects defined naturally by distance from the origin or by direction
often acquire remarkably simple equations of their own.

\subsection*{Circles}

Consider the equation $r = a$, where $a$ is a positive constant. Every
point satisfying this equation lies exactly $a$ units from the origin, so
the set of all such points is a circle of radius $a$ centered at the
origin. A circle that requires $x^2+y^2 = a^2$ in Cartesian coordinates
becomes simply $r = a$ in polar coordinates, a dramatic simplification
that is one of the principal reasons for introducing the system.

\begin{example}
Sketch $r=4$. Every point lies four units from the origin, so the graph
is a circle of radius four centered at the pole.
\end{example}

\subsection*{Lines Through the Origin}

Now consider $\theta = c$, where $c$ is a constant angle. Every point
satisfying this equation lies in the same direction from the origin; the
distance may vary, but the angle remains fixed, so the graph is a line
through the origin making angle $c$ with the positive $x$-axis.

\begin{example}
The equation $\theta = \pi/4$ describes the line through the origin
inclined at $45^\circ$.
\end{example}

Together, $r=a$ and $\theta=c$ illustrate a recurring theme of coordinate
systems: some objects become simpler while others become more
complicated, and the usefulness of a coordinate system depends on whether
its notion of simplicity aligns with the geometry under investigation.

\section{Periodic Functions as Curves}

One of the most interesting consequences of polar coordinates appears
when trigonometric functions are used to define the radius. In
Chapter~6, sine and cosine were studied as periodic functions whose
outputs changed as an angle varied. In polar coordinates, those outputs
can be interpreted as distances from the origin, so as $\theta$ changes,
the radius changes as well, producing a curve in the plane; the process
transforms a periodic function into a geometric figure.

Consider $r = a\cos(n\theta)$ or $r = a\sin(n\theta)$. These equations
generate a family of curves known as \emph{rose curves}: the repeated
oscillation of the sine or cosine function causes the graph to move
alternately toward and away from the origin, producing a pattern of
petals arranged around the pole, with the number of petals depending on
$n$ and the constant $a$ controlling their size. The periodicity that
once produced repeating waves along a horizontal axis now produces
repeating geometric structures distributed around a circle.

\begin{example}
Consider $r = 4\cos(2\theta)$. As $\theta$ varies, the cosine function
oscillates between $-1$ and $1$, so the radius varies between $-4$ and
$4$. At $\theta = 0$, $r = 4$, so the graph begins four units to the
right of the origin. At $\theta = \pi/4$, $r = 0$, so the curve passes
through the pole. At $\theta = \pi/2$, $r = -4$; a negative radius places
the point four units in the \emph{opposite} direction, creating a petal
above the origin. As $\theta$ continues to increase, the same pattern
repeats, and the result is a four-petaled rose centered at the pole.
\end{example}

The rose curve provides a striking example of how periodicity and
geometry interact: a simple oscillation in the radius becomes a symmetric
two-dimensional pattern, and the repetition formerly visible in the graph
of a function over time is now visible as rotational symmetry in space.

\subsection*{Cardioids}

A second important family of polar curves is given by equations of the
form $r = a(1+\cos\theta)$ or $r = a(1+\sin\theta)$, called
\emph{cardioids}, from the Greek word for ``heart,'' because of their
characteristic shape. Unlike a rose curve, a cardioid does not repeatedly
pass through the origin; instead the radius grows and shrinks smoothly,
touching the origin at a single cusp.

\begin{example}
Consider $r = 2(1+\cos\theta)$. At $\theta = 0$, the radius reaches its
maximum value, $r = 4$; at $\theta = \pi$, $r = 0$, so the curve passes
through the pole. As $\theta$ moves between these extremes, the radius
varies continuously, producing the familiar cardioid shape.
\end{example}

Cardioids appear naturally in optics, acoustics, and signal processing;
many microphone pickup patterns, for example, are approximately
cardioid, capturing sound more strongly from one direction than another.

\subsection*{Spirals}

Not all polar curves are periodic. Consider $r = a\theta$: as the angle
increases, the distance from the origin increases as well, so the point
continually rotates while simultaneously moving outward, tracing a curve
called an \emph{Archimedean spiral}.

\begin{example}
The equation $r = \theta$ describes a spiral beginning at the origin. At
$\theta=0$, $r=0$; at $\theta=\pi$, $r=\pi$; at $\theta=2\pi$, $r=2\pi$.
Each complete revolution increases the distance from the origin by the
same amount, causing the curve to wind steadily outward.
\end{example}

Unlike circles, roses, and cardioids, spirals do not close upon
themselves; they continue indefinitely, illustrating that polar
coordinates can naturally describe both periodic and non-periodic
geometric behavior.

\begin{sidebar}{Writing Systems and Alternative Representations}
In Arabic script, a single letter often appears in several different
forms depending on its position within a word: an isolated form used on
its own, and distinct initial, medial, and final forms used according to
which neighboring letters it joins. The letter commonly transliterated
\emph{`ayn}, for instance, has four such shapes, one for each of these
positions, all of which are read as occurrences of the same underlying
letter despite looking visibly different on the page. Polar coordinates
exhibit a directly analogous phenomenon: the pairs $(r,\theta)$,
$(r,\theta+2\pi)$, and $(-r,\theta+\pi)$ may all represent the same
location, even though the coordinate pairs themselves are different. In
both cases, the representation changes according to context while the
underlying object it denotes remains the same; what looks like a
difference in description is not necessarily a difference in the thing
being described.
\end{sidebar}

\section{Applications}

Polar coordinates are particularly useful whenever distance and direction
are more natural than horizontal and vertical displacement. In
navigation, positions are often specified by a bearing and a range: a
lighthouse may be reported as lying ten kilometers away at a bearing of
$40^\circ$ north of east, information already expressed in polar form
and requiring no conversion. In physics, many systems possess rotational
symmetry; the motion of a planet around a star, the orbit of a satellite,
or the behavior of a particle moving under a central force are often
easier to describe using distance from a center and angular position
than by repeatedly tracking horizontal and vertical coordinates.
Engineers frequently use polar plots to display directional information,
since antenna radiation patterns, microphone sensitivity diagrams, and
radar coverage maps are commonly expressed as functions of angle, making
polar coordinates the natural framework for visualizing them. The
usefulness of polar coordinates arises from the same principle that
motivates every coordinate system: a good representation is one in which
the important structure of a problem becomes simple.

\section{Looking Ahead}

Polar coordinates describe a point using two quantities, $(r,\theta)$:
the first records magnitude, the second records direction. Throughout
this chapter these quantities have been treated as separate numbers
joined together into an ordered pair, yet many of the geometric
operations studied earlier naturally involve both simultaneously.
Rotations alter $\theta$ while preserving $r$; scaling changes $r$ while
preserving $\theta$; the trigonometric functions repeatedly connect the
two. This raises a natural question: must magnitude and direction remain
separate, or can they be combined into a single mathematical object? The
next chapter takes up exactly this question, unifying the scaling
transformation of Chapter~1 and the rotation transformation of Chapter~8
into a single operator whose two parameters are precisely the amplitude
and angle already familiar from this chapter, before Chapter~16 shows
that same combination reappearing as ordinary multiplication in an
extended number system.

\section{Exercises}

\subsection*{Converting Coordinates}
\begin{exercise}
Convert each polar coordinate pair to rectangular coordinates: (a)
$\left(4,\dfrac{\pi}{6}\right)$; (b) $\left(6,\dfrac{3\pi}{4}\right)$; (c)
$(5,\pi)$.
\end{exercise}

\begin{exercise}
Convert each rectangular coordinate pair to polar coordinates, giving an
exact value of $\theta$ whenever possible: (a) $(3,3\sqrt3)$; (b)
$(-4,4)$; (c) $(-5,0)$.
\end{exercise}

\begin{exercise}
Convert $(2,-2\sqrt3)$ to polar coordinates.
\end{exercise}

\begin{exercise}
Convert $(-3,-3)$ to polar coordinates.
\end{exercise}

\subsection*{Quadrants and Angle Determination}
\begin{exercise}
For each point, determine the quadrant in which the angle $\theta$ lies
before computing the polar coordinates: (a) $(-2,5)$; (b) $(-7,-1)$; (c)
$(4,-6)$.
\end{exercise}

\begin{exercise}
Explain why the formula $\theta = \arctan(y/x)$ cannot by itself
determine the correct polar angle.
\end{exercise}

\begin{exercise}
Find two different points having the same value of $y/x$ but lying in
different quadrants, and explain why they require different polar
angles.
\end{exercise}

\subsection*{Equivalent Polar Representations}
\begin{exercise}
Find three different polar representations of $\left(3,\dfrac{\pi}{4}\right)$.
\end{exercise}

\begin{exercise}
Show that $\left(5,\dfrac{\pi}{6}\right)$ and
$\left(-5,\dfrac{7\pi}{6}\right)$ represent the same point.
\end{exercise}

\begin{exercise}
Find a positive-radius representation and a negative-radius
representation of the point $\left(2,\dfrac{2\pi}{3}\right)$.
\end{exercise}

\begin{exercise}
Explain why every polar coordinate pair has infinitely many equivalent
representations.
\end{exercise}

\subsection*{Graphing in Polar Coordinates}
\begin{exercise}
Describe the graph of each equation: (a) $r=3$; (b) $r=7$; (c)
$\theta=\dfrac{\pi}{3}$.
\end{exercise}

\begin{exercise}
Sketch the graphs of $r=2$, $r=5$, and $\theta = \dfrac{\pi}{4}$.
\end{exercise}

\begin{exercise}
Explain why the equation $r=a$ always represents a circle centered at
the origin.
\end{exercise}

\begin{exercise}
Explain why the equation $\theta=c$ always represents a line through the
origin.
\end{exercise}

\subsection*{Polar Curves}
\begin{exercise}
For each value of $\theta$, compute the corresponding value of $r$ for
$r = 4\cos(2\theta)$ when $\theta = 0, \dfrac{\pi}{4}, \dfrac{\pi}{2},
\dfrac{3\pi}{4}$.
\end{exercise}

\begin{exercise}
Describe the effect of increasing the constant $a$ in the equation
$r = a\cos(2\theta)$.
\end{exercise}

\begin{exercise}
Sketch a rough graph of $r = 2(1+\cos\theta)$.
\end{exercise}

\begin{exercise}
Compute several points on the spiral $r=\theta$ for $\theta = 0,
\dfrac{\pi}{2}, \pi, \dfrac{3\pi}{2}, 2\pi$.
\end{exercise}

\begin{exercise}
Explain why the spiral $r=\theta$ does not close upon itself, while a
rose curve does.
\end{exercise}

\subsection*{Applications}
\begin{exercise}
A radar station detects an aircraft at distance $120$ km and bearing
$35^\circ$ north of east. Express the aircraft's position in polar
coordinates.
\end{exercise}

\begin{exercise}
A ship is located at $\left(80,\dfrac{\pi}{6}\right)$ in polar
coordinates relative to a lighthouse. Convert the position to
rectangular coordinates.
\end{exercise}

\begin{exercise}
Explain why polar coordinates are often preferable to rectangular
coordinates when describing circular motion.
\end{exercise}

\begin{exercise}
Research one application of polar plots in science or engineering, and
write a short paragraph describing how polar coordinates are used.
\end{exercise}

\subsection*{Conceptual and Exploratory Problems}
\begin{exercise}
Compare the point $(3,4)$ in rectangular coordinates with the same
point represented in polar coordinates. What information is emphasized
by each description?
\end{exercise}

\begin{exercise}
Describe the relationship between the conversion formulas $x=r\cos\theta$,
$y=r\sin\theta$ and the unit-circle definitions of sine and cosine.
\end{exercise}

\begin{exercise}
Why do circles centered at the origin become simpler in polar
coordinates than in rectangular coordinates?
\end{exercise}

\begin{exercise}
The graph of $r = a\cos(n\theta)$ is generated by a periodic function.
Explain how periodicity in $\theta$ produces repeated geometric
structure in the plane.
\end{exercise}

\begin{exercise}
A point may have many different polar representations but only one
rectangular representation. Explain why this difference occurs.
\end{exercise}

\subsection*{Challenge Problems}
\begin{exercise}
Show algebraically that $(r,\theta) = (-r,\theta+\pi)$, using the
conversion formulas to prove the result.
\end{exercise}

\begin{exercise}
Suppose $(r_1,\theta_1)$ and $(r_2,\theta_2)$ represent the same point
and both radii are positive. Prove that $\theta_2 - \theta_1 = 2\pi k$
for some integer $k$.
\end{exercise}

\begin{exercise}
Find all polar representations of the point $(0,0)$, and explain why the
origin is exceptional in the polar coordinate system.
\end{exercise}

\begin{exercise}
Polar coordinates describe a point using distance and direction. Explain
why this makes them a natural precursor to the amplitude-and-angle
representation developed in the next chapter.
\end{exercise}
